\documentclass[aspectratio=169]{beamer}
% SGSSS Beamer Preamble — shared across all lectures
% Brand colours from SGSSS logo

\usetheme{metropolis}

% --- Colour definitions ---
\definecolor{sgsssMagenta}{HTML}{A3217A}
\definecolor{sgsssGreen}{HTML}{2D6A3F}
\definecolor{sgsssBlack}{HTML}{1D1D1B}
\definecolor{sgsssWhite}{HTML}{FFFFFF}
\definecolor{sgsssLightGrey}{HTML}{F5F5F5}

% --- Metropolis colour overrides ---
\setbeamercolor{palette primary}{bg=sgsssMagenta, fg=sgsssWhite}
\setbeamercolor{title separator}{fg=sgsssGreen}
\setbeamercolor{progress bar}{fg=sgsssGreen, bg=sgsssLightGrey}
\setbeamercolor{frametitle}{bg=sgsssMagenta, fg=sgsssWhite}
\setbeamercolor{alerted text}{fg=sgsssMagenta}
\setbeamercolor{example text}{fg=sgsssGreen}
\setbeamercolor{block title}{bg=sgsssMagenta, fg=sgsssWhite}
\setbeamercolor{block body}{bg=sgsssLightGrey, fg=sgsssBlack}
\setbeamercolor{block title example}{bg=sgsssGreen, fg=sgsssWhite}
\setbeamercolor{block body example}{bg=sgsssLightGrey, fg=sgsssBlack}
\setbeamercolor{block title alerted}{bg=sgsssMagenta, fg=sgsssWhite}
\setbeamercolor{block body alerted}{bg=sgsssLightGrey, fg=sgsssBlack}
\setbeamercolor{normal text}{fg=sgsssBlack}

% --- Fonts ---
\usepackage[T1]{fontenc}
\usepackage[sfdefault]{FiraSans}
\usepackage{FiraMono}
\setbeamerfont{frametitle}{size=\large}

% --- Packages ---
\usepackage{graphicx}
\usepackage{bookmark}
\hypersetup{
    bookmarks=true,
    bookmarksopen=true,
    bookmarksdepth=2,
    pdfpagemode=UseOutlines,
}
% --- PDF bookmarks for each frame title ---
\makeatletter
\apptocmd{\beamer@@frametitle}{%
  \only<1>{\bookmark[page=\the\c@page,level=1]{#1}}%
}{}{}
\makeatother
\usepackage{array}
\usepackage{booktabs}
\usepackage{minted}
\usepackage{fontawesome5}

% --- TikZ ---
\usepackage{tikz}
\usetikzlibrary{arrows.meta, positioning, shapes.geometric, trees, calc, decorations.pathmorphing, decorations.pathreplacing, fit, backgrounds}

% Reusable TikZ styles
\tikzset{
    server node/.style={rectangle, draw=sgsssGreen, fill=sgsssGreen!10, thick, minimum width=2.5cm, minimum height=1.2cm, align=center, font=\small},
    client node/.style={rectangle, draw=sgsssMagenta, fill=sgsssMagenta!10, thick, minimum width=2.5cm, minimum height=1.2cm, align=center, font=\small},
    api node/.style={rectangle, draw=sgsssGreen, fill=sgsssGreen!15, thick, rounded corners, minimum width=2.5cm, minimum height=1.2cm, align=center, font=\small},
    data node/.style={rectangle, draw=sgsssBlack, fill=sgsssLightGrey, thick, minimum width=2cm, minimum height=1cm, align=center, font=\small},
    arrow style/.style={-{Stealth[length=3mm]}, thick, sgsssBlack},
    html tag/.style={rectangle, draw=sgsssMagenta, fill=sgsssMagenta!8, rounded corners=2pt, font=\ttfamily\small, inner sep=4pt},
    json key/.style={rectangle, draw=sgsssGreen, fill=sgsssGreen!8, rounded corners=2pt, font=\ttfamily\small, inner sep=4pt},
}

% --- Minted configuration ---
\setminted{
    fontsize=\footnotesize,
    frame=leftline,
    framesep=2mm,
    baselinestretch=1.1,
    bgcolor=sgsssLightGrey,
    breaklines,
    linenos=false,
}
\setminted[python]{style=friendly}
\setminted[r]{style=friendly}

% --- Convenience commands ---
\newcommand{\pythoncode}[1]{\mintinline{python}{#1}}
\newcommand{\rcode}[1]{\mintinline{r}{#1}}
\newcommand{\httpstatus}[2]{\texttt{#1} \textit{#2}}

% --- Footer (page numbers only; logo on title slide only) ---
\setbeamertemplate{frame footer}{%
    \insertframenumber/\inserttotalframenumber%
}

% --- Metadata ---
\author{Dr Diarmuid McDonnell}
\institute{Braw Data Ltd \and Gradel Institute of Charity, University of Oxford}
\date{27 March 2026}


\title{Exploring Text: Descriptive Approaches}
\subtitle{SGSSS Training Course}

\begin{document}

% -------------------------------------------------------
% Slide 1: Section title
% -------------------------------------------------------
{
\setbeamertemplate{frame footer}{}
\begin{frame}[plain]
    \begin{center}
        \vspace{1.5cm}
        {\Large\bfseries\color{sgsssMagenta} Exploring Text}\\[0.3cm]
        {\large Descriptive Approaches}\\[0.8cm]
        {\normalsize Dr Diarmuid McDonnell}\\[0.1cm]
        {\small Braw Data Ltd \(\cdot\) Gradel Institute of Charity, University of Oxford}\\[0.3cm]
        {\small 27 March 2026}
    \end{center}
\end{frame}
}

% -------------------------------------------------------
% Slide 2: Why explore before modelling?
% -------------------------------------------------------
\begin{frame}{Why Explore Before Modelling?}
    \begin{alertblock}{Key idea}
        Just as we compute descriptive statistics before running regressions, we should \textbf{explore text} before applying complex models.
    \end{alertblock}

    \begin{itemize}\setlength\itemsep{0pt}
        \item \faIcon{search} \textbf{Understand your data} --- identify patterns, outliers, and data-quality issues early
        \item \faIcon{bug} \textbf{Catch preprocessing errors} --- missing documents, encoding problems, unexpected languages
        \item \faIcon{lightbulb} \textbf{Generate hypotheses} --- exploratory analysis reveals structure that informs modelling choices
    \end{itemize}

    \vspace{0.2cm}
    {\small Descriptive text analysis is valuable in its own right --- not merely a stepping stone to fancier methods.}
\end{frame}

% -------------------------------------------------------
% Slide 3: Word frequencies
% -------------------------------------------------------
\begin{frame}{Word Frequencies}
    \small
    \begin{columns}[T]
        \begin{column}{0.52\textwidth}
            \textbf{What they show:}
            \begin{itemize}\setlength\itemsep{0pt}
                \item Which terms dominate a corpus
                \item Differences across groups or time periods
                \item Raw counts or relative frequencies
            \end{itemize}

            \vspace{0.2cm}
            \textbf{\color{sgsssGreen} Zipf's law:} a small number of words account for most tokens; most words are rare.

            \vspace{0.2cm}
            \begin{alertblock}{Limitations}
                \footnotesize
                High-frequency words are often \textbf{stopwords}. Stemming and lemmatisation choices affect counts.
            \end{alertblock}
        \end{column}
        \begin{column}{0.44\textwidth}
            \begin{exampleblock}{Example: Top 5 terms}
                \footnotesize
                \begin{tabular}{@{}lr@{}}
                    \toprule
                    \textbf{Term} & \textbf{Frequency} \\
                    \midrule
                    government  & 1{,}842 \\
                    member      & 1{,}536 \\
                    minister    & 1{,}204 \\
                    people      &   987 \\
                    policy      &   763 \\
                    \bottomrule
                \end{tabular}
            \end{exampleblock}
            {\scriptsize After stopword removal; UK parliamentary debates.}
        \end{column}
    \end{columns}
\end{frame}

% -------------------------------------------------------
% Slide 4: Word clouds
% -------------------------------------------------------
\begin{frame}{Word Clouds}
    \begin{columns}[T]
        \begin{column}{0.50\textwidth}
            \textbf{\color{sgsssGreen} Strengths}
            \begin{itemize}
                \item \faIcon{eye} Intuitive and visually engaging
                \item Quick overview of dominant terms
                \item Useful for presentations to non-technical audiences
            \end{itemize}
        \end{column}
        \begin{column}{0.50\textwidth}
            \textbf{\color{sgsssMagenta} Limitations}
            \begin{itemize}
                \item \faIcon{exclamation-triangle} No information about \textbf{context}
                \item Word size is hard to compare precisely
                \item Layout is arbitrary --- can mislead
                \item Not reproducible across runs (random placement)
            \end{itemize}
        \end{column}
    \end{columns}

    \vspace{0.4cm}
    \begin{alertblock}{Rule of thumb}
        Word clouds are for \textbf{communication}, not \textbf{analysis}. Always supplement with frequency tables or bar charts.
    \end{alertblock}
\end{frame}

% -------------------------------------------------------
% Slide 5: Keywords-in-context (KWIC)
% -------------------------------------------------------
\begin{frame}[fragile]{Keywords-in-Context (KWIC)}
    \begin{block}{Definition}
        Concordance lines showing a \textbf{target word} with its surrounding context, revealing \textit{how} terms are used --- not just how often.
    \end{block}

    \vspace{0.1cm}
    \begin{exampleblock}{Example: ``immigration'' in parliamentary debates}
        \footnotesize
        \begin{verbatim}
  ...controls on  immigration  are essential to...
  ...debate about  immigration  has become too...
  ...the impact of immigration  on public services...
  ...reduce illegal immigration  through cooperation...\end{verbatim}
    \end{exampleblock}

    {\small \faIcon{info-circle} KWIC reveals framing, sentiment, and collocational patterns that frequency counts alone miss.}
\end{frame}

% -------------------------------------------------------
% Slide 6: Collocations and n-grams
% -------------------------------------------------------
\begin{frame}{Collocations and N-grams}
    \small
    \begin{alertblock}{Key idea}
        Words that \textbf{co-occur more than expected by chance} carry meaning beyond individual terms.
    \end{alertblock}

    \vspace{-0.1cm}
    \begin{description}[\textbf{Trigrams}]\setlength\itemsep{0pt}
        \item[\textbf{Bigrams}] Two-word sequences: \textit{climate change}, \textit{prime minister}
        \item[\textbf{Trigrams}] Three-word sequences: \textit{cost of living}, \textit{United Nations Security}
    \end{description}

    \vspace{-0.1cm}
    \textbf{Identifying meaningful collocations:}
    \begin{itemize}\setlength\itemsep{0pt}
        \item Raw co-occurrence counts (biased towards frequent words)
        \item \textbf{Pointwise Mutual Information (PMI)} --- how much more often words co-occur than expected under independence
        \item Chi-squared and log-likelihood tests
    \end{itemize}
\end{frame}

% -------------------------------------------------------
% Slide 7: Lexical diversity
% -------------------------------------------------------
\begin{frame}{Lexical Diversity}
    \small
    \begin{block}{Type--Token Ratio (TTR)}
        \(\text{TTR} = \frac{\text{Number of unique types}}{\text{Total number of tokens}}\)
        \quad --- a higher TTR indicates a \textbf{richer, more varied vocabulary}.
    \end{block}

    \textbf{Applications in social science:}
    \begin{itemize}\setlength\itemsep{0pt}
        \item Comparing \textbf{vocabulary richness} across political speakers
        \item Tracking linguistic complexity over time
        \item Assessing readability of policy documents
    \end{itemize}

    \begin{alertblock}{Caveat}
        \footnotesize
        TTR is sensitive to \textbf{document length} --- longer texts tend to have lower TTR. Use standardised measures (e.g., MATTR, moving-average TTR) when comparing texts of different lengths.
    \end{alertblock}
\end{frame}

% -------------------------------------------------------
% Slide 8: Next up
% -------------------------------------------------------
\begin{frame}{Next Up}
    \begin{center}
        {\Large \alert{Practical 2: Descriptive Text Analysis}}\\[0.5cm]
        \begin{itemize}
            \item Compute and visualise word frequencies
            \item Generate KWIC concordance lines
            \item Identify collocations and n-grams
            \item Measure lexical diversity across speakers
        \end{itemize}
        \vspace{0.5cm}
        {\normalsize Open the \textbf{Practical 2} notebook in Google Colab\\(Python or R --- your choice!)}
    \end{center}
\end{frame}

\end{document}
